% ===== Opening: title, hook, the book and the series, today's plan (~5 min) =====

\webinartitle{5}{Գլուխ 6. Անհատականություն. դժվար թեստ}%
  {\chip{masks}{թեստեր} \chip{angry}{զայրույթ} \chip{rocket}{մոտիվացիա} \chip{laughing}{հումոր}}%
  {confused}{Անձը չի գտնվել}{{crystalball}{joker}{volcano}}

% Hook: the online quizzes everybody has taken
\begin{frame}{Ո՞վ ես դու ըստ թեստի}
  \begin{columns}[c]
    \begin{column}{0.56\textwidth}
      \begin{popup}{թեստ.exe}
        {\blackfont\large Իմացիր քո իսկական անհատականությունը}\par\smallskip
        \begin{itemize}
          \item<2-> Ո՞ր պիցցան ես դու \emoji{pizza}
          \item<3-> Հոգվարթսի ո՞ր տնից ես \emoji{tophat}
          \item<4-> Քո 4 տառը. INTJ, ENFP, ISFJ...
          \item<5-> «Ա» տիպի՞ ես, թե՞ «Բ»
        \end{itemize}
        \onslide<2->{{\scriptsize\color{plum!60!cream}12 հարց · 2 րոպե · 100\% գիտական}}
      \end{popup}
    \end{column}
    \begin{column}{0.38\textwidth}
      \centering
      \mascot[2.8cm]{thinking}\par\medskip
      \onslide<6->{%
      \begin{chatbox}
        Անհատականության թեստ անցե՞լ ես։ Գրիր արդյունքդ (օր.՝ \textbf{INFP} կամ «Գրիֆինդոր»)
      \end{chatbox}}
    \end{column}
  \end{columns}
\end{frame}

% The book and where this chapter sits in the series
\begin{frame}{Գիրքը և շարքը}
  \begin{columns}[T]
    \begin{column}{0.5\textwidth}
      \begin{card}
        {\blackfont\Large The Idiot Brain} {\small (2016)}\par
        «Ապուշ ուղեղը». նեյրոգիտնական \textbf{Դին Բըրնեթը} բացատրում է, թե ինչու է ուղեղն այդքան տպավորիչ ու միաժամանակ այդքան անկատար\par\medskip
        \begin{itemize}
          \item Ամեն վեբինար՝ գրքի մեկ գլուխ
          \item Գլուխները իրարից անկախ են. կարելի է միանալ ցանկացածից
          \item Ամեն պնդման տակ գիտական աղբյուր կա
        \end{itemize}
      \end{card}
    \end{column}
    \begin{column}{0.44\textwidth}
      {\blackfont\small Գրքի 8 գլուխը}\par\smallskip
      {\color{plum!55!cream}\small
      Ներածություն \quad 1. Մարմին\\
      2. Հիշողություն \quad 3. Վախ\\
      4. Խելք \quad 5. Ընկալում}\par\smallskip
      \chip[brainpink]{masks}{6. Անհատականություն}\quad{\small\blackfont $\leftarrow$ այսօր}\par\smallskip
      {\color{plum!55!cream}\small
      7. Ուրիշ մարդիկ\\
      8. Երբ ուղեղը «փչանում» է}\par\medskip
      \begin{card}[left=1.3mm, right=1.3mm, colback=lavender!35!cardwhite]
        Այսօր՝ \textbf{ով ենք մենք}. բնավորություն, զայրույթ, մոտիվացիա ու հումոր
      \end{card}
    \end{column}
  \end{columns}
\end{frame}

% Plan for today
\begin{frame}{Այսօր}
  \begin{columns}[T]
    \begin{column}{0.19\textwidth}
      \begin{card}[left=1.3mm, right=1.3mm, colback=butter!45!cream]
        \centering\bigemoji[0.9cm]{railway}\par\smallskip
        {\blackfont Սկիզբ}\par\smallskip
        \raggedright Մետաղյա ձող գլխի միջով
      \end{card}
    \end{column}
    \begin{column}{0.19\textwidth}
      \begin{card}[left=1.3mm, right=1.3mm, colback=lavender!45!cream]
        \centering\bigemoji[0.9cm]{masks}\par\smallskip
        {\blackfont Թեստեր}\par\smallskip
        \raggedright Ինչո՞վ է «Մեծ հնգյակը» տարբերվում MBTI-ից
      \end{card}
    \end{column}
    \begin{column}{0.19\textwidth}
      \begin{card}[left=1.3mm, right=1.3mm, colback=brainpink!40!cream]
        \centering\bigemoji[0.9cm]{angry}\par\smallskip
        {\blackfont Զայրույթ}\par\smallskip
        \raggedright Ինչու է այն երբեմն օգտակար
      \end{card}
    \end{column}
    \begin{column}{0.19\textwidth}
      \begin{card}[left=1.3mm, right=1.3mm, colback=mint!45!cream]
        \centering\bigemoji[0.9cm]{rocket}\par\smallskip
        {\blackfont Մոտիվացիա}\par\smallskip
        \raggedright Ինչն է մեզ իրականում շարժում
      \end{card}
    \end{column}
    \begin{column}{0.19\textwidth}
      \begin{card}[left=1.3mm, right=1.3mm, colback=peach!45!cream]
        \centering\bigemoji[0.9cm]{laughing}\par\smallskip
        {\blackfont Հումոր}\par\smallskip
        \raggedright Ինչու է դժվար ծիծաղը չափել
      \end{card}
    \end{column}
  \end{columns}
  \vspace{4mm}
  \centering
  {\small \emoji{coffee}\ Մեջտեղում՝ 5 րոպե ընդմիջում \qquad \emoji{speech}\ Հարցերը գրիր չաթում ցանկացած պահի}
\end{frame}
